\documentstyle[boldmath,semcolor,semlayer,portrait,rotate,epsf,german,fancybox]{seminar}
\input{unilogo}
\input{amssym.def}
\input{amssym}
\def\dummy#1{}
\overlaystrue
\def\overlay#1{{}}

\slidewidth9.5in
\slideheight6.7in
%\def\slidetopmargin{\iflandscape.2in\else0.8in\fi}
\def\slidetopmargin{\iflandscape.3in\else0.8in\fi}
\def\slidebottommargin{\iflandscape1.0in\else0.4in\fi}
\def\slideleftmargin{\iflandscape.2in\else.5in\fi}
\def\sliderightmargin{\iflandscape1.0in\else.8in\fi}

\slidesmag{4}
\centerslidesfalse
\slideframewidth0.5pt
%\slideframe{none}

%\portraitonly
%\landscapeonly
%\def\leftsliderotation#1{\rotate[l]{#1}}
%\sliderotation{left}
\rotateheaderstrue

%%%%%%%%%%%%%%%%%%%%%%%%%%%%%%%%%%%%%%%%%%%%%%
%% Makros
%%%%%%%%%%%%%%%%%%%%%%%%%%%%%%%%%%%%%%%%%%%%%%
\def\bra#1{\left<#1\right|}
\def\ket#1{\left|#1\right>}
\newcommand\scalar[2]{\left<#1|#2\right>} 
\newcommand{\scalars}[3]{\left<#1|#2|#3\right>}
\newcommand{\betrag}[1]{\vert #1 \vert}  
\newcommand{\unten}[1]{\lfloor #1 \rfloor} 
\newcommand{\oben}[1]{\lceil #1 \rceil} \def\dmin{\mathop{\rm d}_{\rm min}\nolimits}
\newcommand{\diag}[1]{\mbox{ diag}\left(#1\right)} 
\def\dH{\mathop{\rm d}_H\nolimits}
\def\wH{\mathop{\rm wgt}_H\nolimits}

\def\mathsymfont#1{\ifmmode{\mathchoice{%
    \mbox{\the\textfont0\ignorespaces#1}}{%
    \mbox{\the\textfont0\ignorespaces#1}}{%
    \mbox{\the\scriptfont0\ignorespaces#1}}{%
    \mbox{\the\scriptscriptfont0\ignorespaces#1}}}%
  \else{{\rm\ignorespaces#1}}\fi}
\def\C{{\cal C}}
\def\Cdual{\C^{\perp}}
\def\Czdual{\C_2^{\perp}}
\def\Csum{\sum_{\bm{c}\in\C}}
\def\Cdualsum{\sum_{\bm{c}\in\Cdual}}
\def\F{\mathsymfont{I\kern-.25em F}}
\def\ds{\displaystyle}
\newfont{\bsf}{cmssbx10 scaled 1400}
\newfont{\lbsf}{cmssbx10 scaled 1800}
\newfont{\emsf}{cmssi10}
\newfont{\trm}{cmr10}  % \ten\roman
\newfont{\twrm}{cmr12} % \twelve\roman
\newfont{\brm}{cmbx10 scaled 1400} % \bold\roman
\def\R{\mbox{\trm I\kern-.2em R}}
%\def\bm#1{\mbox{\boldmath$#1$}}
\newcommand{\heading}[1]{%
  \begin{center}
    \bsf
    \shadowbox{#1}%
  \end{center}
  \vspace{1ex minus 1ex}%
}
\newbox{\logobox}
\setbox\logobox=\hbox{\raise 0.17\baselineskip\hbox{\epsfbox{logo.eps}}}
\def\logo{\usebox{\logobox}}
\newpagestyle{MH}{\hss\hbox to \hsize{PhysComp96\hss\thepage}\hss}%
{\hss\hbox to \hsize{${\cal IAKS}$\ \logo\ UniKA\hfil
M.~Grassl}\hss}

\newpagestyle{nohead}{}{\hss\hbox to \hsize{\tiny${\cal IAKS}$\ \logo\
UniKA\hfil \thepage/PhysComp96}\hss}

\newpagestyle{nohead_beth}{}{\hss\hbox to \hsize{\tiny${\cal IAKS}$\ \logo\
UniKA\hfil}\hss}

%%%%%%%%%%%%%%%%%%%%%%%%%%%%%%%%%%%%%%%%%%%%%%

\begin{document}
\frenchspacing
\sf

%%\pagestyle{nohead}
\pagestyle{nohead_beth}
%%\pagestyle{empty}

%%%%%%%%%%%%%%%%%%%%%%%%%%%%%%%%%%%%%%%%%%%%%%
%
% Grovers Suchalgorithmus 
%
%%%%%%%%%%%%%%%%%%%%%%%%%%%%%%%%%%%%%%%%%%%%%%

\begin{slide}\sf
\heading{Grovers Suchalgorithmus}
\fbox{Problemstellung}
\begin{itemize}
  \item  Geg.: $C:X_N \to \left\{ 0,1\right\}$, wobei $X_N=\left\{0,1,\ldots ,N-1\right\}$
  \item  Ges.: ein $i\in X_N$ mit $C(i)=1$
\end{itemize}
\vfill
Die Anzahl der L"osungen $\betrag{\{i\mid C(i)=1\}} = t$ sei bekannt.
\vfill
Annahme: $C$ als black-box gegeben.
\vfill
{\bf Klassisch} (deterministisch oder probabilistisch) ${\cal O}(N)$ Schritte  \\
{\bf Quantencomputer} ${\cal O}(\sqrt{N})$ Schritte
\vfill
\end{slide}

%%%%%%%%%%%%%%%%%%%%%%%%%%%%%%%%%%%%%%%%%%%%%%
%
% Grovers Suchalgorithmus 
%
%%%%%%%%%%%%%%%%%%%%%%%%%%%%%%%%%%%%%%%%%%%%%%

\begin{slide*}\sf
\heading{Unit"are Transformationen}
\begin{center}
\fbox{die bedingte Phasendrehung $S_C$}
\end{center}
   $$ S_C \ket{i} = \left\{ 
                      \begin{array}{rc}
                            -\ket{i} & \mbox{ falls $C\left( i\right) = 1$} \\ 
                             \ket{i} & \mbox{ sonst }
                      \end{array}
                   \right.
  $$
\begin{center} \fbox{Herstellung der Gleichverteilung} \end{center}
$$  W\ket{0} = \frac{1}{\sqrt{N}} \sum_{i=0}^{N-1} \ket{i}                 
             = \frac{1}{\sqrt{N}}\pmatrix{1      & *     & \ldots & *      \cr 
                                          \vdots & \vdots& \ddots & \vdots \cr
                                          1      & *     & \ldots & *}
$$
\begin{center}\fbox{Diffusionstransformation $D$}\end{center}
\begin{eqnarray*}
  D & = & W \mbox{diag}\{1,-1,\ldots,-1\} W^{\dagger} \\
    & = &  \pmatrix{ -1+\frac{2}{N} & \frac{2}{N}   & \ldots      & \frac{2}{N}   \cr 
                               \frac{2}{N}   & \ddots        & \ddots      & \vdots        \cr 
                               \vdots        & \ddots        & \ddots      & \frac{2}{N}   \cr 
                               \frac{2}{N}   & \ldots        & \frac{2}{N} & -1+\frac{2}{N} \cr}
\end{eqnarray*}
\begin{center}\fbox{Grovertransformation $G$}\end{center}
  $$ G_C = D S_C. $$ 
\end{slide*}

%%%%%%%%%%%%%%%%%%%%%%%%%%%%%%%%%%%%%%%%%%%%%%
%
% Grovers Suchalgorithmus 
%
%%%%%%%%%%%%%%%%%%%%%%%%%%%%%%%%%%%%%%%%%%%%%%

\begin{slide}\sf
\heading{Grovers Suchalgorithmus}
  \fbox{Algorithmus}
  \begin{enumerate}
    \item Pr"apariere den Zustand $\ket{0}$.
    \item Stelle Gleichverteilung durch 
          $ W\ket{0} = \frac{1}{\sqrt{N}}\sum_{i=0}^{N-1} \ket{i} $ her.
    \item Wende Grovertransformation $G_C$ insgesamt ${\cal O}(N/t)$ mal an.
    \item Messung des Registers liefert $i_0$.
    \item Pr"ufe, ob $C(i_0)=1$, falls nicht, dann Neustart.
  \end{enumerate} 
\end{slide}

%%%%%%%%%%%%%%%%%%%%%%%%%%%%%%%%%%%%%%%%%%%%%%
%
% Grovers Suchalgorithmus 
%
%%%%%%%%%%%%%%%%%%%%%%%%%%%%%%%%%%%%%%%%%%%%%%

\begin{slide}\sf
\heading{Diffusionstransformation}
   Die Diffusionstransformation $D$ l"a\3t sich wie folgt zerlegen 
\begin{eqnarray*}
   D & = & \pmatrix{ -1+\frac{2}{N} & \frac{2}{N}   & \ldots      & \frac{2}{N}   \cr 
                      \frac{2}{N}   & \ddots        & \ddots      & \vdots        \cr 
                      \vdots        & \ddots        & \ddots      & \frac{2}{N}   \cr 
                      \frac{2}{N}   & \ldots        & \frac{2}{N} & -1+\frac{2}{N} \cr} \\
        & = & D=-I+2P,
\end{eqnarray*} 
wobei $I=\sum_{i=0}^{N-1} \ket{i}\bra{i}$ und $P= \frac{1}{N} \sum_{i,j=0}^{N-1} \ket{i}\bra{j}$.
\end{slide}

\begin{slide}\sf
\heading{Inversion am Durchschnittswert}
   $$P\ket{\Psi} = \left( \frac{1}{N} \sum_{j=0}^{N-1}\scalar{j}{\Psi} \right) \sum_{i=0}^{N-1} \ket{i}$$

   Die Diffusionstransformation $D$ realisiert eine Inversion am Durchschnittswert $\alpha$, denn
   \begin{eqnarray*} 
     D\ket{\Psi} & = & \left( -I+2P\right) \ket{\Psi} \\ 
                 & = & -\ket{\Psi} + 2P\ket{\Psi} . 
   \end{eqnarray*} 
   F"ur die $i$-te Komponente dieses Zustands gilt 
   \begin{eqnarray*} 
     \scalars{i}{D}{\Psi} & = & -\scalar{i}{\Psi} + 2\alpha \\  
                          & = & \left( \alpha + \left( \alpha - \scalar{i}{\Psi}\right)\right) . 
   \end{eqnarray*} 
\end{slide}


%%%%%%%%%%%%%%%%%%%%%%%%%%%%%%%%%%%%%%%%%%%%%%
%
% Grovers Suchalgorithmus 
% 
%%%%%%%%%%%%%%%%%%%%%%%%%%%%%%%%%%%%%%%%%%%%%%
\begin{slide*}\sf
\heading{Amplituden}
\begin{tabular}{|p{0.45\textwidth}|p{0.45\textwidth}|}
\hline
"Uberlagerung \newline 
aller $i\in X_N$ \newline
Hadamardtransf.&
\begin{picture}(100,20)(0,25)
\put(0,0){\vector(1,0){90}}
\put(0,0){\vector(0,1){20}}\thicklines
\multiput(5,0)(5,0){16}{\line(0,1){10}}
\end{picture}\\
\hline&\\[-3.5ex]
\hline
Phasendrehung \newline 
f"ur $C(i)=1$ \newline
&
\begin{picture}(100,15)(0,15)
\put(0,0){\vector(1,0){90}}
\put(0,0){\vector(0,1){20}}\thicklines
\multiput(5,0)(5,0){7}{\line(0,1){10}}
\put(40,0){\line(0,-1){10}}
\multiput(45,0)(5,0){3}{\line(0,1){10}}
\put(60,0){\line(0,-1){10}}
\multiput(65,0)(5,0){4}{\line(0,1){10}}
\end{picture}\\
\hline
Grovertransf. $G_C$\newline Inversion an $\alpha$&
%%\newline{}&
\begin{picture}(100,20)(0,20)
\put(0,5){\vector(1,0){90}}
\put(0,5){\vector(0,1){20}}\thicklines
\multiput(5,5)(5,0){7}{\line(0,1){8}}
\put(40,5){\line(0,1){12}}
\multiput(45,5)(5,0){3}{\line(0,1){8}}
\put(60,5){\line(0,1){12}}
\multiput(65,5)(5,0){4}{\line(0,1){8}}
\end{picture}\\
\hline&\\[-3.5ex]
\hline
\centerline{$\vdots$}&\centerline{$\vdots$}\\
\hline&\\[-3.5ex]
\hline
nach einigen Iterationen\newline{} &
\begin{picture}(100,25)(0,10)
\put(0,5){\vector(1,0){90}}
\put(0,5){\vector(0,1){20}}\thicklines
\multiput(5,5)(5,0){7}{\line(0,1){2}}
\put(40,5){\line(0,1){20}}
\multiput(45,5)(5,0){3}{\line(0,1){2}}
\put(60,5){\line(0,1){20}}
\multiput(65,5)(5,0){4}{\line(0,1){2}}
\end{picture}\\
\hline
\end{tabular}

\bigskip
Wird ein Fixpunkt erreicht?
\end{slide*}

%%%%%%%%%%%%%%%%%%%%%%%%%%%%%%%%%%%%%%%%%%%%%%
%
% Grovers Suchalgorithmus
% 
%%%%%%%%%%%%%%%%%%%%%%%%%%%%%%%%%%%%%%%%%%%%%%
\begin{slide*}\sf
\heading{Amplituden}
   Die geschlossenen Formeln f"ur die Amplituden der L"osungen $k_j$ und
   der Nicht-L"osungen $l_j$ sind
   \begin{eqnarray*} 
      k_{j} & = & \frac{1}{\sqrt{t}}\sin\left(\left( 2j + 1\right)\Theta\right) \\ 
      l_{j} & = & \frac{1}{\sqrt{N-t}}\cos\left(\left( 2j + 1\right)\Theta\right),  
   \end{eqnarray*} 
   wobei der Winkel $\Theta$ gleich $\sin^2\Theta =\frac{t}{N}$ ist. \\

   Erfolgswahrscheinlichkeit am gr"o\3ten, wenn $\betrag{l_{\tilde m}}^2$ minimal.
   $$ l_{\tilde m} = 0 \iff \tilde m = \left( \pi - 2\Theta \right) / 4\Theta $$ 
   Da die Anzahl der Iterationen $m$ eine nat"urliche Zahl sein mu\3, wird $m$ folgenderma\3en definiert 
   $$ m = \unten{\pi / 4\Theta}.$$ 
   Grovers Algorithmus hat eine Laufzeit von ${\cal O}\left(\sqrt{\frac{N}{t}}\right)$.
\end{slide*}

%%%%%%%%%%%%%%%%%%%%%%%%%%%%%%%%%%%%%%%%
%%
%% Median Absch"atzung
%%
%%%%%%%%%%%%%%%%%%%%%%%%%%%%%%%%%%%%%%%%

\begin{slide}\sf
\heading{Medianabsch"atzung}
\fbox{Problemstellung}
\begin{itemize}
  \item Geg.: ein unsortiertes Array $V[i]$ mit $N$ verschiedenen Elementen aus einer total geordneten Menge $({\cal V},\le )$
  \item Ziel: Finde den Median, das $\oben{N/2}$-gr"o\3te Element in $V$ bzgl. $\le$.
\end{itemize}
\end{slide}           

%%%%%%%%%%%%%%%%%%%%%%%%%%%%%%%%%%%%%%%%
%%
%% Median Absch"atzung
%%
%%%%%%%%%%%%%%%%%%%%%%%%%%%%%%%%%%%%%%%%

\begin{slide}\sf
\heading{Medianabsch"atzung}
   Partitionierung des Arrays $V$ durch einen Schwellenwertindex $\mu\in{\cal V}$ in 
   $$ A=\{V[i] | V[i]<\mu\}  \mbox{ und in } B=\{V[i] | V[i]\ge\mu\}.$$ 

    Sei $\epsilon$ so definiert, da\3 $\betrag{A}=(N/2)(1+\epsilon)$ und 
    $\betrag{B}=(N/2)(1-\epsilon)$. 
 
    Dabei gelte $0.1\epsilon_0 < \betrag{\epsilon} < \epsilon_0$ f"ur ein festes $\epsilon_0$, das kleiner als $0.1$ ist.

\end{slide}

%%%%%%%%%%%%%%%%%%%%%%%%%%%%%%%%%%%%%%%%
%%
%% Median Absch"atzung
%%
%%%%%%%%%%%%%%%%%%%%%%%%%%%%%%%%%%%%%%%%

\begin{slide}\sf
\heading{Medianabsch"atzung}
  \fbox{Aufgabe}

  Bestimme $\betrag{\epsilon}$ mit einer Genauigkeit $\Delta$.

  F"ur alle Sch"atzwerte $\widetilde\epsilon$ f"ur $\betrag{\epsilon}$ soll mit hoher
  Wahrscheinlichkeit $\widetilde{\epsilon}(1-\Delta) < \betrag{\epsilon} < \widetilde{\epsilon}(1+\Delta)$ gelten.

  L"osung dieses Problems durch Sch"atzung der Anzahl der Zust"ande $\ket{i}$ mit $C(i)=0$, wobei
  $$ C(i) = \cases{ 1 & falls $V[i] \ge\mu$ \cr
                    0 & falls $V[i] < \mu$ }. $$
  QC-Algorithmus ben"otigt ${\cal O}(\frac{1}{\epsilon_0 \Delta^2})$ Proben.

  Klassische Algorithmen ben"otigen $\Omega(\frac{1}{\epsilon_0^2 \Delta^2})$ Proben.
  
\end{slide}

%%%%%%%%%%%%%%%%%%%%%%%%%%%%%%%%%%%%%%%%
%%
%% Median Absch"atzung
%%
%%%%%%%%%%%%%%%%%%%%%%%%%%%%%%%%%%%%%%%%

\begin{slide}\sf   
\heading{Algorithmus}
         \begin{enumerate}
         \item $\nu={\cal O}(\frac{1}{\Delta^2})$ Messwerte ermitteln.
         \begin{enumerate}
           \item Gleichverteilung $ \ket{\Psi_0} = \sum_{i=0}^{N-1} \ket{i} $
           \item bedingte Phasendrehung $S_{C,\frac{\pi}{2}}$ 
           \item Schiebetransformation $T$ 
           \item Wiederhole die folgenden Schritte $m={\cal O}(\frac{1}{\epsilon_0})$ mal.
             \begin{enumerate}
               \item bedingte Phasendrehung $S_{\neg C,\pi}$ 
               \item Diffusionstransformation $D$ 
               \item bedingte Phasendrehung $S_{C,\pi}$
               \item Diffusionstransformation $D$ 
             \end{enumerate}
           \item Messung des Registers und Speicherung der Messwerte
         \end{enumerate}
         \item Sch"atzung von $\betrag{\epsilon}$ durch statistische Auswertung der Messwerte
         \end{enumerate}
\end{slide}

%%%%%%%%%%%%%%%%%%%%%%%%%%%%%%%%%%%%%%%%
%%
%% Median Absch"atzung
%%
%%%%%%%%%%%%%%%%%%%%%%%%%%%%%%%%%%%%%%%%

\begin{slide}\sf
\heading{Funktionsweise des Algorithmus}
\fbox{1. Teil : Messwerte ermitteln}
  $$ 
  \begin{array}{rlcc}
    \mbox{1.(a)} & \mbox{Herstellung der Gleichverteilung }            & k=1         & l=1 \cr
    \mbox{1.(b)} & \mbox{bedingte Phasendrehung } S_{C,\frac{\pi}{2}}  & k=1         & l=i \cr
    \mbox{1.(c)} & \mbox{Schiebetransformation } T=-iI+(1+i)P          & k=\epsilon  & l=(1+\epsilon)+i
  \end{array}
  $$

\end{slide}

%%%%%%%%%%%%%%%%%%%%%%%%%%%%%%%%%%%%%%%%
%%
%% Median Absch"atzung
%%
%%%%%%%%%%%%%%%%%%%%%%%%%%%%%%%%%%%%%%%%

\begin{slide}\sf
\heading{Funktionsweise des Algorithmus}
\fbox{1.(d) Lineare Verst"arkung der Betr"age der $k$-Amplituden}

Die geschlossenen Formeln f"ur die Amplituden sind
\begin{eqnarray*}
  k_j & = & \epsilon\cos(j\Theta) + \gamma(1+\epsilon+i)\sin(j\Theta) \\
  l_j & = & -\frac{\epsilon}{\gamma}\sin{j\Theta}+(1+\epsilon+i)\cos(j\Theta),
\end{eqnarray*}
wobei $\cos(\Theta) = 1 - 2\epsilon^2$ und $\gamma^2 = \frac{2\epsilon-2\epsilon^2}{2\epsilon+2\epsilon^2}$.

Es gilt f"ur den Betrag der $k$-Amplituden
$$ \betrag{k_j} \approx 2\sqrt{2}\betrag{\epsilon}j, $$
da $\epsilon$ klein ist.

\end{slide}

%%%%%%%%%%%%%%%%%%%%%%%%%%%%%%%%%%%%%%%%
%%
%% Median Absch"atzung
%%
%%%%%%%%%%%%%%%%%%%%%%%%%%%%%%%%%%%%%%%%

\begin{slide}\sf
\heading{Funktionsweise des Algorithmus}
\fbox{Anzahl der Iterationen}

Man kann ein $m < \frac{1}{40\epsilon_0}$ so w"ahlen, so da\3 
$\betrag{k_m} > \frac{\betrag{\epsilon}}{20\epsilon_0}$ gilt.
\vfill
\fbox{2. Teil: Statistische Auswertung der Messwerte}

Die Wahrscheinlichkeit, einen Zustand $\ket{i}$ mit $C(i)=0$ 
zu messen, betr"agt $f(\epsilon)=\frac{1}{2}(\betrag{k_m }^2(1+\epsilon))$, 
da es $(N/2)(1+\epsilon)$ solcher Zust"ande gibt. 
\vfill
Der Wert von $\betrag{\epsilon}$ kann mit einem maximal erwarteten Fehler 
von ${\cal O}(\Delta\epsilon_0)$ aus dem Sch"atzwert $f$ f"ur 
$f(\epsilon)$ gewonnen werden.
\vfill
\end{slide}

%%%%%%%%%%%%%%%%%%%%%%%%%%%%%%%%%%%%%%%%
%%
%% Median Absch"atzung
%%
%%%%%%%%%%%%%%%%%%%%%%%%%%%%%%%%%%%%%%%%

\begin{slide*}\sf
\heading{Medianbestimmung}
Durch Variation des Schwellwerts $\mu$ kann $\epsilon$ bestimmt werden.
\begin{verbatim}
median(min, max)
{
  lower = min;
  upper = max;
  while (upper - lower > 1)
  {
    mu = (upper + lower) / 2;
    // liefert epsilon mit Vorzeichen
    est = epsilon_est(mu, 0.1, 0.1); 
    if (est > 0) upper = mu;
    else         lower = mu;
  }
  return mu;
}
\end{verbatim}
\end{slide*}

\end{document}


